\subsection{不变量定理}

设 g 是李代数，\(\rho\) 是 g 在向量空间 M 上的一个表示。对于 g 的每个单表示类 \(\delta\)，令 \(M_{\delta}\) 为 M 的 \(\delta\) 型分量。M 的不变元子空间 \(M_{0}\) 正是 \(M_{\delta_{0}}\)，其中 \(\delta_{0}\) 表示 g 在一维空间上的零表示类。

引理 4. 设 \(\rho, \sigma, \tau\) 是 \(\mathfrak{g}\) 在向量空间 \(\mathbf{M}, \mathbf{N}, \mathbf{P}\) 上的表示。假设给定一个 K-双线性映射 \((m, n) \mapsto m \cdot n\)，它把 \(\mathbf{M} \times \mathbf{N}\) 映入 \(\mathbf{P}\)，并满足

\[
(\rho (x) m). n + m. (\sigma (x) n) = \tau (x) (m. n)
\]

对所有 \(m \in M\)、\(n \in N\) 和 \(x \in g\) 都成立。

\begin{enumerate}
\def\labelenumi{(\alph{enumi})}
\item
  若 \(m_0 \in \mathbf{M}_0\)，则映射 \(n \mapsto m_0\).\(n\) 是 \(\mathfrak{g}\)-模同态。
\item
  若 \(n \in N_{\delta}\)，则 \(m_{0}.n \in P_{\delta}\)。
\item
  若 \(\mathbf{M}\) 是一个（不必结合的）代数，且 \(\rho(x)\) 是 \(\mathbf{M}\) 的导子，同时 \(\mathbf{M}_0\) 是 \(\mathbf{M}\) 的子代数，则每个 \(\mathbf{M}_{\delta}\) 都是 \(\mathbf{M}_0\) 的左、右模。
\end{enumerate}

对于 \(m_{0} \in M_{0}, n \in N\) 和 \(x \in g,\)，有：

\[
\tau (x) \left(m _ {0}. n\right) = m _ {0}. (\sigma (x) n),
\]

因此得到（a）。断言（b）由（a）推出（《代数》，第八章，§ 3，第 4 号，命题 10）。若 \(\mathbf{N} = \mathbf{P} = \mathbf{M}\) 且 \(\sigma = \tau = \rho\)，则断言（b）给出断言（c），作为特殊情形。

引理 5. 进一步设 \(\sigma\) 和 \(\tau\) 是半单的，因而 N（分别为 P）是 \(N_{\delta}\)（分别为 \(P_{\delta}\)）的直和。对于所有 \(n \in N\)（分别对于 \(p \in P\)），令 \(n^{\sharp}\)（分别为 \(p^{\sharp}\)）表示它在 \(N_0\)（分别为 \(P_0\)）中的分量。设 \(m_0 \in M_0\)。则对所有 \(n \in N\)，有 \((m_0.n)^{\sharp} = m_0.n^{\sharp}\)。

由线性性，只需考虑 \(n \in \mathrm{N}_{\delta}\) 的情形。若 \(\delta \neq \delta_0\)，则 \(n^{\sharp} = 0\)，且由引理 4，\(m_0.n \in \mathrm{P}_{\delta}\)，所以 \((m_0.n)^{\sharp} = 0 = m_0.n^{\sharp}\)。若 \(\delta = \delta_0\)，则 \(n^{\sharp} = n\)，且由引理 4，\(m_0.n \in \mathrm{P}_0\)，所以 \((m_0.n)^{\sharp} = m_0.n = m_0.n^{\sharp}\)。

定理 6. 设 \(\mathfrak{g}\) 是李代数，V 是 \(\mathfrak{g}\)-模，是 K 上的有限维半单模；S 是 V 的对称代数，\(x_{\mathrm{S}}\) 是 S 上延拓 \(x_{\mathrm{V}}\) 的导子（从而 \(x \mapsto x_{\mathrm{S}}\) 是 \(\mathfrak{g}\) 在 S 上的一个表示）。

\begin{enumerate}
\def\labelenumi{(\alph{enumi})}
\item
  \(\mathbf{S}_0\) 的不变量代数 \(\mathbf{S}\) 由有限个元素生成。
\item
  对于每个单表示类 \(\delta\)（它是 \(\mathfrak{g}\) 的、且在 \(\mathbf{K}\) 上有限维的单表示类），令 \(\mathrm{S}_{\delta}\) 为 \(\mathrm{S}\) 的 \(\delta\) 型分量。则 \(\mathrm{S}_{\delta}\) 是有限生成的 \(\mathrm{S}_{0}\)-模。
\end{enumerate}

令 \(\overline{\mathbf{S}}\subset \mathbf{S}\) 为 S 中无常数项元素构成的理想。令 I 为由 \(\mathrm{S}_0\cap \overline{\mathrm{S}}\) 生成的 S 的理想，令 \((s_1,s_2,\ldots ,s_p)\) 为 I 的一组有限生成元（《交换代数》，第三章，§ 3）。可以假设 \(s_t\) 属于 \(\mathbf{S}_0 \cap \bar{\mathbf{S}}\) 且是齐次的（\(x_{\mathrm{S}}\) 保持次数，因而每个 \(\mathbf{S}_0\) 都是分次子模）。令 \(\mathbf{S}_1\) 为 \(\mathbf{S}\) 的子代数，由 1 和 \(s_t\) 生成。于是 \(\mathbf{S}_1 \subset \mathbf{S}_0\)。我们证明 \(\mathbf{S}_1 = \mathbf{S}_0\)。为此，我们证明：每个齐次元素 \(s\) 若属于 \(\mathbf{S}_0\)，便属于 \(\mathbf{S}_1\)；证明按 \(n\)（即 \(s\) 的次数）归纳。若 \(n = 0\)，断言显然成立。设 \(n > 0\)，并假设当次数小于 n（即 \(< n\)）时断言已成立。由于 \(s \in I\)，有 \(s = \sum_{i=1}^{p} s_i s_i'\)，其中 \(s_i'\) 是 S 中的元素，并可设为齐次的，且 \(\deg(s_i') = \deg(s) - \deg(s_i) < n\)。

\[
s = s ^ {\sharp} = \sum_ {i = 1} ^ {p} (s _ {i} s _ {i} ^ {\prime}) ^ {\sharp} = \sum_ {i = 1} ^ {p} s _ {i} {s _ {i} ^ {\prime}} ^ {\sharp}.
\]

\(s_i^\sharp\) 是 \(S_0\) 中次数 \(< n\) 的齐次元素（因为每个 \(S_\delta\) 都是分次子模），因此根据归纳假设它们属于 \(S_1\)。所以 \(s \in S_1\)，从而证明（a）。

现在考察 g 的一个 \(\delta\) 类的单表示，它作用在有限维空间 M 上。令 \(L = \mathcal{L}_{\mathbb{K}}(M, S)\)。对于所有 \(s \in S\) 和所有 \(f \in L\)，令 sf 为 L 中由 \((sf)(m) = s.f(m)\) 定义的元素（\(m \in M\)）；于是 L 上定义了 S-模结构。由于 M 在 K 上有限维，显然 L 是有限生成的 S-模，因而是诺特 S-模（因为环 S 是诺特的）。另一方面，L 具有规范的 g-模结构。对每个整数 \(n \geqslant 0\)，令 \(S^{n}\) 为 S 中次数 n 的齐次元素集合；于是，g-模 \(\mathcal{L}_{\mathbb{K}}(M, S^{n})\) 是半单的（第 5 号，定理 4 的推论 3），因而 g-模 L 是半单的。此外，对于 \(s \in S\)、\(f \in L\)、\(x \in g\) 和 \(m \in M\)，

\[
\begin{array}{r l} (x _ {\mathrm{L}} (s f)) (m) & = x _ {\mathrm{S}} ((s f) (m)) - (s f) (x _ {\mathrm{M}} m) \\ & = x _ {\mathrm{S}} (s. f (m)) - s. f (x _ {\mathrm{M}} m) \\ & = (x _ {\mathrm{S}} s). f (m) + s. (x _ {\mathrm{S}} f (m)) - s. f (x _ {\mathrm{M}} m) \\ & = ((x _ {\mathrm{S}} s) f) (m) + (s (x _ {\mathrm{L}} f)) (m) \end{array}
\]

从而 \(x_{\mathrm{L}}(sf) = (x_{\mathrm{S}}s)f + s(x_{\mathrm{L}}f)\)。因此可应用引理 5。

L 的不变元子集 \(\mathbf{L}_0\) 正是 \(\mathbf{L}\) 中从 \(g\)-模 \(\mathbf{M}\) 到 \(g\)-模 \(\mathbf{S}\) 的同态集合。因此，若 \(\phi\) 表示从 \(\mathrm{M} \otimes_{\mathrm{K}} \mathrm{L}\) 到 \(\mathrm{S}\) 的规范同态，则 \(\phi(\mathrm{M} \otimes_{\mathrm{K}} \mathrm{L}_0) = \mathrm{S}_{\delta}\)。由于 \(\phi\) 显然是 S-模同态，只需证明 \(\mathbf{L}_0\) 是有限生成的 \(\mathrm{S}_0\)-模。令 \(\mathbf{J}\) 为 \(\mathbf{L}\) 的由 \(\mathbf{L}_0\) 生成的 S-子模。由于 \(\mathbf{L}\) 是诺特 S-模，存在一个有限序列 \((f_1, \ldots, f_q)\)，其元素属于 \(\mathbf{L}_0\)，并生成 S-模 \(\mathbf{J}\)。令 \(\mathbf{L}_1\) 为 \(\mathrm{S}_0\)-模，由 \(f_{1},\ldots ,f_{q}\) 生成。于是 \(\mathrm{L}_1\subset \mathrm{L}_0\)。另一方面，若 \(f\in \mathrm{L}_0\)，则 \(f = \sum_{i = 1}^{q}s_{i},f_{i}\)，其中 \(s_i\in S\) 对所有 \(i\) 成立；因此，根据引理 5（沿用其记号）：

\[
f = f ^ {\sharp} = \left(\sum_ {i = 1} ^ {q} s _ {i} f _ {i}\right) ^ {\sharp} = \sum_ {i = 1} ^ {q} s _ {i} ^ {\sharp} f _ {i} \in L _ {1}.
\]

因此 \(L_{0} = L_{1}\)，所以 \(L_{0}\) 是有限生成的 \(S_{0}\)-模。

